\section{Introduction}
An image communication system with a hard per-image symbol budget must decide which information to retain before channel transmission. The decision depends on the image-dependent source-code length, channel protection, and the reconstruction attainable from the retained representation. A generative receiver can supply missing information, but its output must remain constrained by the information actually delivered. The resulting design trades pixel fidelity, perceptual similarity, feature consistency, and processing cost.

Digital image coding with forward error correction and learned joint source--channel coding (JSCC) offer different approaches to this tradeoff~\cite{deepjscc,swinjscc}. Progressive transmission and truncation of embedded descriptions have a long history~\cite{chande}. Recent token-based methods use large generative models: Ada-TokenCom combines prefix entropy coding, tail generation, and joint prefix/MCS selection under an average resource constraint~\cite{adatokencom}; TONIC studies task-oriented token communication~\cite{tonic}. ARPC constructs progressive image descriptions with an autoregressive prior~\cite{arpc}. Visual autoregressive modeling (VAR) provides a multiscale residual representation and a prior that predicts each scale from complete coarser maps~\cite{var}.

For this next-scale representation, restricting transmission to complete maps gives large endpoint steps. Meanwhile, the terminated entropy-code length of a given endpoint varies across images. A next whole scale may exceed the frame even when the preceding whole endpoint leaves substantial room. Legal within-scale endpoints address this mismatch by retaining part of the boundary map. The next-scale factorization lets both ends construct matching CDFs from complete coarse maps, while the extra endpoints make this representation fit a hard frame more closely.

We present NeST-Com, which connects within-scale endpoints, actual-length fitting, and prefix-constrained reconstruction in one fixed-budget communication design. The receiver uses the existing known-token overwrite operation of VAR while updating both the accumulated residual latent and the context for finer scales. Our contributions are the endpoint/frame-fitting design and its evaluation through shared-model comparisons, controlled granularity diagnostics, and six-system experiments. On common500, quarter-scale endpoints improve PSNR by 0.32--0.70~dB over complete-scale endpoints. External comparisons expose complementary pixel and perceptual behavior rather than a uniform ranking.
